\begin{examplebox}
	\textbf{\textcolor{CExample}{例 1（基础）}}
	\textbf{\textcolor{CExample}{题目：}} 已知 $\vec{e}_1,\vec{e}_2$ 不共线，$\vec{a}=2\vec{e}_1-\vec{e}_2$，$\vec{b}=\vec{e}_1+3\vec{e}_2$，求 $\vec{a}+\vec{b}$ 和 $2\vec{a}-\vec{b}$ 用 $\vec{e}_1,\vec{e}_2$ 表示。\textbf{\textcolor{CExample}{解题步骤：}}
	\begin{description}[style=nextline, leftmargin=2.8em, labelsep=0.5em, itemsep=0.45em]
		\item[\textbf{第一步（代入加法式）：}]
			写 $\vec{a}+\vec{b} = (2\vec{e}_1-\vec{e}_2) + (\vec{e}_1+3\vec{e}_2)$。
			\par\textbf{理由：}直接代入向量加法。
		\item[\textbf{第二步（合并加法式）：}]
			$\vec{a}+\vec{b} = 3\vec{e}_1+2\vec{e}_2$。
			\par\textbf{理由：}合并同类项。
		\item[\textbf{第三步（差倍运算）：}]
			\[
				\begin{aligned}
					2\vec{a}-\vec{b}&= 2(2\vec{e}_1-\vec{e}_2) - (\vec{e}_1+3\vec{e}_2) \\&= (4\vec{e}_1-2\vec{e}_2) - \vec{e}_1 - 3\vec{e}_2 \\&= (4-1)\vec{e}_1 + (-2-3)\vec{e}_2 \\&= 3\vec{e}_1 - 5\vec{e}_2.
			\end{aligned}
			\]
			\par\textbf{理由：}数乘分配律与合并。
	\end{description}
	\textbf{\textcolor{CExample}{关键结论：}}
	\textbf{$\vec{a}+\vec{b}=3\vec{e}_1+2\vec{e}_2$，$2\vec{a}-\vec{b}=3\vec{e}_1-5\vec{e}_2$。}

	\medskip
	\textbf{\textcolor{CExample}{例 2（稍变形）}}
	\textbf{\textcolor{CExample}{题目：}} 在 $\triangle ABC$ 中，$D$ 是 $BC$ 上一点，且 $BD=2DC$，用 $\overrightarrow{AB},\overrightarrow{AC}$ 表示 $\overrightarrow{AD}$。\textbf{\textcolor{CExample}{解题步骤：}}
	\begin{description}[style=nextline, leftmargin=2.8em, labelsep=0.5em, itemsep=0.45em]
		\item[\textbf{第一步（选取基底）：}]
			以 $\overrightarrow{AB},\overrightarrow{AC}$ 为基底，它们不共线。
			\par\textbf{理由：}三角形两边不共线，可作为基底。
		\item[\textbf{第二步（表示$\overrightarrow{BD}$）：}]
			由 $BD=2DC$ 得 $\overrightarrow{BD}=\frac{2}{3}\overrightarrow{BC}$。又 $\overrightarrow{BC}=\overrightarrow{AC}-\overrightarrow{AB}$，故 $\overrightarrow{BD}=\frac{2}{3}(\overrightarrow{AC}-\overrightarrow{AB})$。
			\par\textbf{理由：}共线比例与向量减法。
		\item[\textbf{第三步（表达$\overrightarrow{AD}$）：}]
			\[
				\begin{aligned}
					\overrightarrow{AD}&= \overrightarrow{AB} + \overrightarrow{BD} \\&= \overrightarrow{AB} + \frac{2}{3}(\overrightarrow{AC}-\overrightarrow{AB}) \\&= \frac{1}{3}\overrightarrow{AB} + \frac{2}{3}\overrightarrow{AC}.
			\end{aligned}
			\]
			\par\textbf{理由：}向量加法与数乘运算，合并同类项。
	\end{description}
	\textbf{\textcolor{CExample}{关键结论：}}
	\textbf{$\overrightarrow{AD} = \frac{1}{3}\overrightarrow{AB} + \frac{2}{3}\overrightarrow{AC}$。}

	\medskip
	\textbf{\textcolor{CExample}{例 3（综合应用）}}
	\textbf{\textcolor{CExample}{题目：}} 设 $\vec{e}_1,\vec{e}_2$ 不共线，$\vec{a}=2\vec{e}_1+3\vec{e}_2$，$\vec{b}=4\vec{e}_1-5\vec{e}_2$，$\vec{c}=10\vec{e}_1-7\vec{e}_2$，判断 $\vec{a},\vec{b}$ 是否不共线，并求 $\vec{c}$ 用 $\vec{a},\vec{b}$ 表示的表达式。\textbf{\textcolor{CExample}{解题步骤：}}
	\begin{description}[style=nextline, leftmargin=2.8em, labelsep=0.5em, itemsep=0.45em]
		\item[\textbf{第一步（基底共线判定）：}]
			由于 $\frac{2}{4}\neq\frac{3}{-5}$，故 $\vec{a},\vec{b}$ 不共线，可作为一组基底。
			\par\textbf{理由：}不共线等价于对应系数不成比例。
		\item[\textbf{第二步（设元展开）：}]
			设 $\vec{c}=x\vec{a}+y\vec{b}$，代入得
			\[
				\begin{aligned}
					\vec{c} &= x(2\vec{e}_1+3\vec{e}_2)+y(4\vec{e}_1-5\vec{e}_2) \\&= (2x+4y)\vec{e}_1+(3x-5y)\vec{e}_2.
			\end{aligned}
			\]
			所以 $(2x+4y)\vec{e}_1+(3x-5y)\vec{e}_2=10\vec{e}_1-7\vec{e}_2$。
			\par\textbf{理由：}线性表示假设，利用向量运算展开。
		\item[\textbf{第三步（比较系数求解）：}]
			由分解唯一性，系数对应相等：
			\[
				\begin{cases}
					2x+4y = 10,\\3x-5y = -7.
			\end{cases}
			\]
			解得 $x=1,\ y=2$。故 $\vec{c}=\vec{a}+2\vec{b}$。
			\par\textbf{理由：}平面向量基本定理保证分解唯一，列方程组求解。
	\end{description}
	\textbf{\textcolor{CExample}{关键结论：}}
	\textbf{$\vec{c}=\vec{a}+2\vec{b}$。}
\end{examplebox}
